% Diagram Descriptions for ESP32 Indoor Positioning System
% This file contains detailed descriptions for all diagrams in the report

\section*{Diagram Reference Guide}

\subsection*{1. System Architecture Diagram (system-architecture.tex)}
\textbf{Purpose:} Illustrates the three-phase evolution of the ESP32 IPS architecture.

\textbf{Description:}
This diagram shows the progressive development of the system through three distinct phases:
\begin{itemize}
    \item \textbf{Phase 1 - Hands-Free (Single Target):} Demonstrates unidirectional communication where a single target node broadcasts to four anchor nodes, which forward RSSI measurements to a gateway, ultimately sending aggregated data to a PC for logging and analysis.
    
    \item \textbf{Phase 2 - Multi-Target Sequential:} Shows time-slotted operation where multiple targets (A, B, C) broadcast sequentially in designated time windows (0-100ms, 100-200ms, 200-300ms) to prevent signal collision. Anchors tag RSSI data with target IDs before forwarding to the gateway.
    
    \item \textbf{Phase 3 - All-Nodes Mesh:} Depicts a fully distributed mesh network where all nodes perform bidirectional pairwise RSSI measurements, creating an adjacency map that enables relative positioning without fixed anchor infrastructure.
\end{itemize}

\textbf{Usage:} Include in Chapter 4 (Proposed Solution) to explain system architecture evolution.

\textbf{LaTeX Integration:}
\begin{verbatim}
\begin{figure}[H]
    \centering
    \begin{tikzpicture}[
    node distance=1.5cm and 2cm,
    box/.style={rectangle, draw, thick, minimum width=2cm, minimum height=0.8cm, align=center, font=\small},
    phase/.style={rectangle, draw, dashed, thick, inner sep=0.5cm},
    arrow/.style={->, thick},
    bidir/.style={<->, thick}
]

% Phase 1: Hands-Free (Single Target)
\node[phase, label={[font=\bfseries]above:Phase 1: Hands-Free (Single Target)}] (phase1) at (0,0) {
    \begin{tikzpicture}[node distance=1cm and 1.5cm, box/.style={rectangle, draw, thick, minimum width=2cm, minimum height=0.8cm, align=center, font=\small}, arrow/.style={->, thick}]
        \node[box, fill=blue!20] (t1) {Target\\Node};
        \node[box, fill=green!20, below left=of t1] (a1) {Anchor 1};
        \node[box, fill=green!20, below=of a1] (a2) {Anchor 2};
        \node[box, fill=green!20, below right=of t1] (a3) {Anchor 3};
        \node[box, fill=green!20, below=of a3] (a4) {Anchor 4};
        \node[box, fill=orange!20, below=2.5cm of t1] (gw1) {Gateway};
        \node[box, fill=red!20, below=of gw1] (pc1) {PC/Logger};
        
        \draw[arrow] (t1) -- node[left, font=\tiny] {Broadcast} (a1);
        \draw[arrow] (t1) -- node[right, font=\tiny] {Broadcast} (a3);
        \draw[arrow] (t1) -- (a2);
        \draw[arrow] (t1) -- (a4);
        \draw[arrow] (a1) -- node[left, font=\tiny] {RSSI} (gw1);
        \draw[arrow] (a2) -- (gw1);
        \draw[arrow] (a3) -- node[right, font=\tiny] {RSSI} (gw1);
        \draw[arrow] (a4) -- (gw1);
        \draw[arrow] (gw1) -- node[right, font=\tiny] {Data} (pc1);
    \end{tikzpicture}
};

% Phase 2: Multi-Target Sequential
\node[phase, label={[font=\bfseries]above:Phase 2: Multi-Target Sequential}, right=4cm of phase1] (phase2) {
    \begin{tikzpicture}[node distance=1.2cm and 1.5cm, box/.style={rectangle, draw, thick, minimum width=2cm, minimum height=0.8cm, align=center, font=\small}, arrow/.style={->, thick}]
        \node[box, fill=blue!20] (t2a) {Target A\\Slot 0-100ms};
        \node[box, fill=blue!20, below=0.4cm of t2a] (t2b) {Target B\\Slot 100-200ms};
        \node[box, fill=blue!20, below=0.4cm of t2b] (t2c) {Target C\\Slot 200-300ms};
        \node[box, fill=green!20, below=2.2cm of t2b] (an) {Anchors\\(All)};
        \node[box, fill=orange!20, below=1.2cm of an] (gw2) {Gateway};
        \node[box, fill=red!20, below=1.2cm of gw2] (pc2) {PC/Logger};
        
        % Curved arrows to avoid overlap
        \draw[arrow] (t2a.south) to[out=-90, in=150] (an.north west);
        \draw[arrow] (t2b.south) -- (an.north);
        \draw[arrow] (t2c.south) to[out=-90, in=30] (an.north east);
        \draw[arrow] (an) -- node[right, font=\tiny, align=center] {Tagged\\RSSI} (gw2);
        \draw[arrow] (gw2) -- node[right, font=\tiny, align=center] {Multi-Target\\Data} (pc2);
    \end{tikzpicture}
};

% Phase 3: All-Nodes Mesh
\node[phase, label={[font=\bfseries]above:Phase 3: All-Nodes Mesh}, right=4cm of phase2] (phase3) {
    \begin{tikzpicture}[node distance=2cm and 2cm]
        \node[box, fill=purple!20] (n1) {Node 1};
        \node[box, fill=purple!20, right=of n1] (n2) {Node 2};
        \node[box, fill=purple!20, below=of n1] (n4) {Node 4};
        \node[box, fill=purple!20, below=of n2] (n3) {Node 3};
        \node[box, fill=orange!20, below=1.5cm of n4] (gw3) {Gateway};
        
        \draw[bidir] (n1) -- node[above, font=\tiny] {RSSI} (n2);
        \draw[bidir] (n2) -- node[right, font=\tiny] {RSSI} (n3);
        \draw[bidir] (n3) -- node[below, font=\tiny] {RSSI} (n4);
        \draw[bidir] (n4) -- node[left, font=\tiny] {RSSI} (n1);
        \draw[bidir] (n1) -- (n3);
        \draw[bidir] (n2) -- (n4);
        
        \draw[arrow] (n1) -- (gw3);
        \draw[arrow] (n2) -- (gw3);
        \draw[arrow] (n3) -- (gw3);
        \draw[arrow] (n4) -- (gw3);
    \end{tikzpicture}
};

% Evolution arrows
\draw[->, ultra thick, blue] ([xshift=0.5cm]phase1.east) -- ([xshift=-0.5cm]phase2.west) node[midway, above] {Evolution};
\draw[->, ultra thick, blue] ([xshift=0.5cm]phase2.east) -- ([xshift=-0.5cm]phase3.west) node[midway, above] {Evolution};

\end{tikzpicture}

    \caption{Three-Phase System Architecture Evolution}
    \label{fig:system-architecture}
\end{figure}
\end{verbatim}

\subsection*{2. Data Communication Flow Diagram (data-communication-flow.tex)}
\textbf{Purpose:} Illustrates the sequence of data communication between system components.

\textbf{Description:}
This sequence diagram shows the temporal flow of data through the system:
\begin{enumerate}
    \item Target node broadcasts ESP-NOW packets containing target\_id, timestamp, and message
    \item All anchor nodes (A1-A4) receive the broadcast simultaneously
    \item Each anchor records RSSI value and timestamp
    \item Anchors forward anchor\_data\_t structures (anchor\_id, rssi, timestamp) to gateway
    \item Gateway aggregates data from all anchors
    \item Gateway sends complete dataset (X, Y, RSSI\_A1, RSSI\_A2, RSSI\_A3, RSSI\_A4) to PC
\end{enumerate}

\textbf{Key Data Structures:}
\begin{itemize}
    \item \texttt{target\_packet\_t}: Broadcast packet from target
    \item \texttt{anchor\_data\_t}: RSSI measurement from anchor
    \item Dataset format: CSV with coordinates and RSSI values
\end{itemize}

\textbf{Usage:} Include in Chapter 4 to explain data flow and packet structures.

\subsection*{3. Hardware Layout Diagram (hardware-layout.tex)}
\textbf{Purpose:} Shows the physical testbed configuration with coordinate system.

\textbf{Description:}
This diagram presents the 70×50 cm testbed layout with:
\begin{itemize}
    \item Four anchor nodes at corners: A1(0,0), A2(70,0), A3(0,50), A4(70,50)
    \item Cartesian coordinate system with 10 cm grid spacing
    \item Sample target positions marked on the grid
    \item Dashed lines showing signal paths from targets to anchors
    \item Dimensional annotations for testbed size
    \item Legend distinguishing anchor and target nodes
\end{itemize}

\textbf{Technical Specifications:}
\begin{itemize}
    \item Testbed dimensions: 70 cm × 50 cm
    \item Grid spacing: 10 cm
    \item Antenna height: 10 cm (consistent across all nodes)
    \item Coordinate origin: Bottom-left corner (A1)
\end{itemize}

\textbf{Usage:} Include in Chapter 4 to show experimental setup and Chapter 5 for calibration reference.

\subsection*{4. Data Flow Diagram - Level 0 (dfd-level0.tex)}
\textbf{Purpose:} Provides context-level view of the entire system and its external entities.

\textbf{Description:}
This context diagram shows the ESP32 IPS System as a single process interacting with three external entities:
\begin{itemize}
    \item \textbf{System Administrator:} Provides configuration parameters and receives system status reports
    \item \textbf{Target Nodes:} Send broadcast signals and receive position estimates
    \item \textbf{Data Analyst:} Receives RSSI datasets and position data, provides calibration constants
\end{itemize}

\textbf{Data Flows:}
\begin{itemize}
    \item Configuration Parameters → System
    \item System Status Reports → Administrator
    \item Broadcast Signals → System
    \item Position Estimates → Target Nodes
    \item RSSI Dataset/Position Data → Data Analyst
    \item Calibration Constants → System
\end{itemize}

\textbf{Usage:} Include in Chapter 4 to provide high-level system overview.

\subsection*{5. Data Flow Diagram - Level 1 (dfd-level1.tex)}
\textbf{Purpose:} Detailed decomposition of system processes and data stores.

\textbf{Description:}
This diagram expands the Level 0 context into five detailed processes:

\textbf{Processes:}
\begin{enumerate}
    \item \textbf{1.0 Broadcast Management:} Manages time-slot allocation and ESP-NOW packet transmission
    \item \textbf{2.0 RSSI Measurement:} Anchor nodes receive broadcasts and measure signal strength
    \item \textbf{3.0 Data Aggregation:} Gateway collects and formats RSSI data from all anchors
    \item \textbf{4.0 Position Estimation:} Applies ML model or trilateration to estimate coordinates
    \item \textbf{5.0 Calibration Processing:} Derives path-loss constants from controlled measurements
\end{enumerate}

\textbf{Data Stores:}
\begin{itemize}
    \item \textbf{D1 - RSSI Database:} Stores raw and aggregated RSSI measurements
    \item \textbf{D2 - Position Records:} Contains estimated X,Y coordinates with timestamps
    \item \textbf{D3 - Config Parameters:} Holds system settings and calibration constants
\end{itemize}

\textbf{Usage:} Include in Chapter 4 for detailed system design documentation.

\subsection*{6. Entity-Relationship Diagram (er-diagram.tex)}
\textbf{Purpose:} Defines the database schema for storing system data.

\textbf{Description:}
This ER diagram shows five main entities and their relationships:

\textbf{Entities:}
\begin{itemize}
    \item \textbf{NODE:} ESP32 devices (targets, anchors, gateway) with unique IDs, MAC addresses, and physical locations
    \item \textbf{MEASUREMENT:} Individual RSSI measurements with timestamps and packet IDs
    \item \textbf{POSITION:} Estimated X,Y coordinates with confidence scores and estimation method
    \item \textbf{CALIBRATION:} Path-loss model parameters (A, n) with calibration dates
    \item \textbf{SESSION:} Experimental sessions grouping measurements by phase and time period
\end{itemize}

\textbf{Relationships:}
\begin{itemize}
    \item One NODE transmits many MEASUREMENTs (1:N)
    \item Many MEASUREMENTs generate one POSITION (N:1)
    \item One CALIBRATION is used by many NODEs (1:N)
    \item One SESSION contains many MEASUREMENTs (1:N)
\end{itemize}

\textbf{Usage:} Include in Chapter 4 or Chapter 5 to document data storage structure.

\subsection*{7. Flowchart - Calibration Protocol (flowchart-calibration.tex)}
\textbf{Purpose:} Details the step-by-step calibration procedure.

\textbf{Description:}
This flowchart illustrates the calibration process for deriving path-loss constants:

\textbf{Process Steps:}
\begin{enumerate}
    \item Place single target at known position (X, Y)
    \item Initialize sample counter to 0
    \item Loop: Collect 50-100 RSSI samples per position
    \item Target broadcasts ESP-NOW packets
    \item All anchors measure RSSI values
    \item Record X, Y, RSSI\_A1, RSSI\_A2, RSSI\_A3, RSSI\_A4
    \item Move target to next grid position
    \item Repeat until all positions covered
    \item Aggregate all samples into dataset
    \item Calculate average RSSI for each position
    \item Derive path-loss constants A and n using regression
    \item Output calibration parameters
\end{enumerate}

\textbf{Key Parameters:}
\begin{itemize}
    \item Sample size: 50-100 per position
    \item Grid spacing: 10 cm
    \item Broadcast interval: 100 ms
    \item Static target (no motion during sampling)
\end{itemize}

\textbf{Path-Loss Model:} $d = 10^{(A - RSSI)/(10n)}$

\textbf{Usage:} Include in Chapter 5 (Experimental Results) to document calibration methodology.

\subsection*{8. Flowchart - Position Estimation (flowchart-position-estimation.tex)}
\textbf{Purpose:} Shows the real-time position estimation algorithm.

\textbf{Description:}
This flowchart details the position estimation process with two alternative methods:

\textbf{Main Process:}
\begin{enumerate}
    \item Input: RSSI values from all anchors
    \item Apply RSSI filtering (moving average)
    \item Validate RSSI values
    \item Select estimation method (ML or Trilateration)
\end{enumerate}

\textbf{ML Path:}
\begin{enumerate}
    \item Normalize RSSI using StandardScaler
    \item Apply MLPRegressor model (64, 64 hidden layers)
    \item Output X, Y coordinates
\end{enumerate}

\textbf{Trilateration Path:}
\begin{enumerate}
    \item Convert RSSI to distance using calibrated A, n parameters
    \item Check if distance values are reasonable
    \item Calculate position using trilateration algorithm
    \item Output X, Y coordinates
\end{enumerate}

\textbf{Post-Processing:}
\begin{enumerate}
    \item Validate position within testbed bounds (0-70 cm, 0-50 cm)
    \item Apply boundary correction if needed
    \item Output final position with timestamp
    \item Store in position database
\end{enumerate}

\textbf{Usage:} Include in Chapter 4 or Chapter 5 to explain position estimation algorithms.

\subsection*{9. Flowchart - System Operation (flowchart-system-operation.tex)}
\textbf{Purpose:} Illustrates the complete system operation for Phase 2 (Multi-Target).

\textbf{Description:}
This flowchart shows the operational flow for multi-target sequential positioning:

\textbf{Initialization:}
\begin{enumerate}
    \item Initialize ESP-NOW protocol
    \item Load configuration (time slots, node IDs)
    \item Synchronize all nodes to common time base
\end{enumerate}

\textbf{Main Operation Loop:}
\begin{enumerate}
    \item Check current time slot
    \item If target's time slot: Broadcast ESP-NOW packet, wait for next slot
    \item If anchor node: Receive broadcast, measure RSSI, send data to gateway
    \item Gateway: Aggregate data from all anchors
    \item Wait for all anchor responses (with timeout)
    \item Estimate position using selected algorithm
    \item Output position data to PC/Logger
    \item Store in database
    \item Repeat
\end{enumerate}

\textbf{Time Slot Configuration:}
\begin{itemize}
    \item Target A: 0-100 ms
    \item Target B: 100-200 ms
    \item Target C: 200-300 ms
    \item Cycle repeats every 300 ms
\end{itemize}

\textbf{Usage:} Include in Chapter 4 to explain system operation and timing.

\subsection*{Integration Guidelines}

\textbf{Chapter 4 - Proposed Solution:}
\begin{itemize}
    \item System Architecture Diagram (mandatory)
    \item Data Communication Flow Diagram (mandatory)
    \item Hardware Layout Diagram (mandatory)
    \item DFD Level 0 and Level 1
    \item ER Diagram
    \item Flowchart - System Operation
\end{itemize}

\textbf{Chapter 5 - Experimental Results:}
\begin{itemize}
    \item Flowchart - Calibration Protocol (mandatory)
    \item Flowchart - Position Estimation
    \item Hardware Layout Diagram (reference)
\end{itemize}

\textbf{LaTeX Tips:}
\begin{itemize}
    \item Use \texttt{[H]} placement for exact positioning
    \item Scale diagrams if needed: \texttt{\textbackslash scalebox\{0.8\}\{\textbackslash input\{...\}\}}
    \item Rotate large diagrams: \texttt{\textbackslash rotatebox\{90\}\{\textbackslash input\{...\}\}}
    \item Always include captions and labels for cross-referencing